\section{工程清单：把研究原型变成生产系统}

\subsection{上线前必须回答的十个问题}
在研究原型阶段，系统往往追求尽快看到成功案例；但走向生产时，必须把鲁棒性、成本、可审计性前置。下面这张清单可作为上线前评审模板。

\begin{table}[H]
\centering
\small
\begin{tabular}{p{0.8cm}p{3.9cm}p{5.8cm}}
\toprule
\textbf{\#} & \textbf{问题} & \textbf{通过标准示例} \\
\midrule
1 & 任务边界是否清晰？ & 输入约束、禁止动作、终止条件均有文档化定义 \\
2 & 状态是否可回放？ & 每步 observation/action/result 均可重放与审计 \\
3 & 搜索预算是否可控？ & 设置最大步数、最大分支、超时回退策略 \\
4 & Judge 是否校准？ & 离线对照集上精度达标，且有定期复核机制 \\
5 & 失败能否自动恢复？ & 支持异常检测、回滚、重规划三段式恢复 \\
6 & 人工介入触发是否明确？ & 对支付、隐私、低置信度场景有强制升级策略 \\
7 & 数据回流是否闭环？ & 成功/失败轨迹都能回流并参与下一轮训练 \\
8 & 成本是否可解释？ & 每类任务有 token、时延、外部调用成本统计 \\
9 & 评估是否持续更新？ & benchmark 与真实流量都做持续监控 \\
10 & 安全与合规是否覆盖？ & 敏感操作有最小权限和审计日志 \\
\bottomrule
\end{tabular}
\caption{Web Agent 生产化前的最小工程清单}
\end{table}

\subsection{最容易被低估的三类工程成本}
\begin{enumerate}
    \item \textbf{状态管理成本}：无状态原型很快，带回滚与多分支状态后复杂度陡增；
    \item \textbf{评估成本}：构建高质量对照集、持续校准 Judge 比训练一个 demo 更费时；
    \item \textbf{运维成本}：网页结构变化频繁，解析规则和动作映射需要持续维护。
\end{enumerate}

\begin{knowledgebox}{为何 ``运维'' 会成为主成本}
很多团队在 PoC 阶段只看模型效果，忽视网页变化导致的规则失效。上线后真正消耗人力的，常常是 selector 变更、页面流程改版、反自动化机制升级等持续运维问题。
\end{knowledgebox}

\subsection{从指标驱动到价值驱动}
讲者谈到 ``economically valuable tasks''，这提醒我们：不要只看 benchmark 分数，而要把任务价值纳入优化目标。一个高分但低价值任务，不应占用大量推理预算；反之，高价值任务应允许更高搜索成本和更强人工协同。

\begin{importantbox}{价值分层调度建议}
\begin{itemize}
    \item 低价值任务：轻量策略 + 低预算 + 快速失败；
    \item 中价值任务：中等搜索 + 自动回滚；
    \item 高价值任务：强搜索 + 多重验证 + 人工介入。
\end{itemize}
\end{importantbox}

\subsection{本章小结}
Agent 生产化的核心难点不在 ``能否跑通 demo''，而在 ``能否稳定、可控、可审计地持续运行''。工程清单和价值分层调度是跨过这道门槛的关键抓手。

